\documentclass[11pt,a4paper]{article}

\usepackage[utf8]{inputenc}
\usepackage[T1]{fontenc}
\usepackage[margin=2.5cm]{geometry}
\usepackage{amsmath,amssymb,bm}
\usepackage{booktabs}
\usepackage{array}
\usepackage{enumerate}
\usepackage{xcolor}
\usepackage{hyperref}
\usepackage{tikz}

\hypersetup{colorlinks=true, linkcolor=blue, urlcolor=blue}

\newcommand{\points}[1]{\textbf{(#1\,\%)}}
\newcommand{\subpoints}[1]{\textit{(#1\,\%)}}
\newcommand{\sol}{\noindent\textbf{Solution}}

\title{TMA4268 V2026 Mock Exam 10 --- Solution Proposal\\[0.3em]
\large}
\author{Compiled for Anders Bekkevard\\[0.2em]
\small Companion to \texttt{mock-exam-10.tex} (same directory).}
\date{Mock for: May 18, 2026}

\begin{document}
\maketitle

\noindent\textit{This document is a worked-solution proposal in the style of the official Stefanie/Sara solutions
to the 2024 and 2025 finals. Point values are quoted in the heading of each solved sub-part. Partial-credit
hints are inline. Refer back to \texttt{mock-exam-10.tex} for the problem statements.}

\vspace{0.6em}
\hrule
\vspace{1em}

%=========================================================
\section*{Problem 1 \points{10} --- Fill-in-the-blank concepts}
%=========================================================

\sol{} \emph{(1\,P per blank.)}

\begin{enumerate}[(1)]
  \item \textbf{parametric}
  \item \textbf{prediction}
  \item \textbf{inference}
  \item \textbf{irreducible}
  \item \textbf{double descent}
  \item \textbf{generative}
  \item \textbf{LDA}
  \item \textbf{bagging}
  \item \textbf{boosting}
  \item \textbf{nested cross-validation}
\end{enumerate}

\noindent\emph{Grading: 1\,P per correct blank, no partial credit for ``close'' alternatives. (4) ``Bayes-optimal''
is a hard zero --- Bayes error and irreducible error are not the same object (Bayes error is for
classification, irreducible error is the regression noise variance, although the two concepts are kin).
(7) ``QDA'' is the classic trap; QDA is the version \emph{without} the shared covariance assumption.
(10) Plain ``cross-validation'' is wrong because the question explicitly asks for a procedure that
separates the selection step from the assessment step.}

\vspace{0.6em}
\hrule
\vspace{1em}

%=========================================================
\section*{Problem 2 \points{28} --- Multiple choice, T/F, short numeric}
%=========================================================

\subsection*{a) Bias--variance, noise, and benign overfitting \subpoints{3}}
\sol{} \emph{(0.75\,P per statement.)}
\begin{enumerate}[(i)]
  \item \textbf{True}. With large irreducible noise $\sigma^2$ a flexible model wastes capacity
        fitting noise that no estimator can recover; the bias--variance trade-off tilts toward
        simpler (higher-bias, lower-variance) models. This is the prof's standard ``high-noise
        regime favours less flexibility'' rule.
  \item \textbf{False}. $\sigma^2$ is the variance of the noise $\varepsilon$ in
        $y_0 = f(x_0) + \varepsilon$; it is a property of the data-generating process and does
        not depend on $n$. The estimator-variance term $\mathrm{Var}(\hat f(x_0))$ typically
        shrinks like $1/n$, but the \emph{irreducible} term does not.
  \item \textbf{False}. The expected squared test error has a third additive term, the
        irreducible noise $\sigma^2$, which is the same for both estimators. With equal bias
        and lower variance, $\hat f_A$ has lower expected squared error \emph{at $x_0$}, but
        ``strictly lower'' is too strong when the difference $\mathrm{Var}(\hat f_B) -
        \mathrm{Var}(\hat f_A)$ could be infinitesimal; more importantly the statement reads as
        if it ignored $\sigma^2$. \emph{Accept ``True'' with a clean argument that the noise
        term cancels and the two variance terms differ strictly.}
  \item \textbf{True}. This is the prof's headline benign-overfitting story: the implicit
        regularization of SGD steers an over-parameterized network toward a minimum-norm
        interpolator among the manifold of zero-training-loss solutions, and that bias is one
        of the structural ingredients that lets test error keep falling past the classical
        bias--variance bound.
\end{enumerate}

\emph{Grading: (iii) is the deliberately ambiguous one; full credit for any answer accompanied by
the correct identification that the $\sigma^2$ term is shared and only the variances differ.}

\subsection*{b) Cross-validation: bias, variance, failure modes, one-SE rule \subpoints{4}}
\sol{}
\begin{enumerate}[(i)]
  \item \subpoints{1} The ten fold MSEs sum to $22.00$, so
        $$\mathrm{CV}_{10} = \frac{22.00}{10} = \boxed{2.20}.$$
        Sample standard deviation (divisor $k-1 = 9$):
        $$s = \sqrt{\frac{\sum (v_i - 2.20)^2}{9}} = \sqrt{\frac{0.31}{9}} \approx 0.186,
          \qquad \widehat{\mathrm{SE}}(\mathrm{CV}_{10}) = \frac{s}{\sqrt{10}} \approx \boxed{0.06}.$$
        \emph{Grading: 0.5\,P for CV-MSE, 0.5\,P for SE. Two decimals as requested.}

  \item \subpoints{1} The CV minimum is at $\lambda = 0.10$ with $\mathrm{CV}_{10} = 2.15$.
        The one-SE bound is
        $$2.15 + 0.06 = \boxed{2.21}.$$
        Values within the band: $\lambda = 0.10$ (CV $= 2.15$) and $\lambda = 0.20$ (CV $= 2.18$).
        The one-SE rule picks the \emph{simplest} (largest-$\lambda$) model among those, hence
        $$\boxed{\hat\lambda_{1\mathrm{SE}} = 0.20}.$$
        \emph{Grading: 0.5\,P for the bound, 0.5\,P for the right $\lambda$. Picking $\lambda = 0.10$
        (the CV-min) is a hard miss --- the whole point of the rule is that we step away from the
        minimum toward a simpler model.}

  \item \subpoints{1} (A) \textbf{True}. LOOCV trains on $n-1$ points (nearly the full sample $\Rightarrow$
        low bias as an estimate of the test error of a size-$n$ trained model), but the $n$ training
        sets overlap almost completely, so the $n$ fold residuals are highly correlated and their
        average has higher variance than the $5$- or $10$-fold average.
        (B) \textbf{False}. With temporally autocorrelated data, random folds break the time order
        and let future information leak into the training of a model that is then evaluated on its
        own past --- the standard remedy is time-series CV (rolling-origin / blocked folds), not a
        bigger $k$.
        \emph{Grading: 0.5\,P each.}

  \item \subpoints{1} The per-fold MSEs are \emph{not} independent (the training sets used to
        produce them overlap heavily), so the formula $s/\sqrt{k}$, which assumes independent
        observations, underestimates the true uncertainty of $\mathrm{CV}_K$.
        \emph{Grading: 1\,P. Accept ``the folds share most of their training data, so the per-fold
        MSEs are positively correlated and $s/\sqrt{k}$ is a stand-in not a real SE.''}
\end{enumerate}

\subsection*{c) Bootstrap: pairs vs.\ residuals, percentile CI \subpoints{3}}
\sol{}
\begin{enumerate}[(i)]
  \item \subpoints{1} \textbf{Paired (case-resampling) bootstrap}. Reason: the residual
        bootstrap assumes the model's functional form is correct and that residuals are i.i.d.\
        and exchangeable across $x$; with random $X$ and a possibly mis-specified model the
        paired scheme makes the weaker assumption that the rows themselves are i.i.d., so it
        remains valid under model mis-specification and under heteroscedastic noise.
        \emph{Grading: 0.5\,P for the right scheme, 0.5\,P for the reason.}

  \item \subpoints{1} The percentile $95\%$ CI is
        $$\Bigl[\hat\theta^{*}_{(0.025)},\;\hat\theta^{*}_{(0.975)}\Bigr]
          \;=\; \Bigl[\hat\theta^{*}_{(25)},\;\hat\theta^{*}_{(975)}\Bigr],$$
        i.e.\ the $2.5\%$ and $97.5\%$ empirical quantiles of the $B = 1000$ bootstrap
        replicates (the $25$th and $975$th order statistics).
        \emph{Grading: 1\,P. ``Mean $\pm 1.96\cdot\mathrm{SE}^*$'' is a \emph{different} CI (normal-approximation), worth $0$ here.}

  \item \subpoints{1} (A) \textbf{True}. By definition $\widehat{\mathrm{SE}}_B(\hat\theta) = s(\hat\theta^*_1,\dots,\hat\theta^*_B)$.
        (B) \textbf{True}. The course recommendation was $B \in [1000,\,10000]$ for SE estimates,
        with more for quantile-based CIs.
        (C) \textbf{False}. The bootstrap estimates the \emph{distribution} of $\hat\theta - \theta$
        (or of $\hat\theta^* - \hat\theta$); it does not correct bias automatically. Bias-corrected
        bootstrap CIs exist (BCa) but the plain percentile CI does not adjust for it.
        \emph{Grading: 0.33\,P each.}
\end{enumerate}

\subsection*{d) Backpropagation, mini-batch SGD, learning rate \subpoints{3}}
\sol{} \emph{(0.75\,P per statement.)}
\begin{enumerate}[(i)]
  \item \textbf{True}. Backprop is the chain-rule recipe for computing $\partial L / \partial \theta$
        layer by layer; the optimizer (SGD, Adam, $\ldots$) is a separate step that consumes that
        gradient.
  \item \textbf{True}. $\widehat{\nabla L}_m = \tfrac{1}{m}\sum_{i\in\mathcal B}\nabla L_i$ is an
        unbiased estimator of the full-data gradient with variance $\propto 1/m$; reducing $m$
        cheapens each step at the cost of a noisier gradient (the noise itself is one source of
        the implicit regularization the prof flagged).
  \item \textbf{False}. $\eta = 2$ is typically two orders of magnitude too large for tabular
        feed-forward training; standard ranges are $10^{-4}$ to $10^{-2}$. Too-large $\eta$ causes
        divergence, not faster convergence.
  \item \textbf{True}. $\partial L / \partial \hat f = \hat f - y$ when $L = \tfrac{1}{2}(\hat f - y)^2$,
        and this residual is exactly the quantity that seeds the backward pass.
\end{enumerate}

\subsection*{e) Collinearity, eigenvalues, dummy coding \subpoints{3}}
\sol{}
\begin{enumerate}[(i)]
  \item \subpoints{1} \textbf{True}. Perfect collinearity means a non-trivial null vector of
        $\mathbf{X}$ (here $\mathbf{X}_2 - 3\mathbf{X}_1 + 2\cdot\mathbf{1} = 0$), so
        $\mathbf{X}^\top \mathbf{X}$ is singular and the OLS normal equations have infinitely many
        solutions; only the predicted value $\hat\beta_1 X_1 + \hat\beta_2 X_2$ (and hence
        $\hat{\mathbf{y}}$) is uniquely identified.
  \item \subpoints{1} \textbf{False}. The collinearity diagnosis is correct, but the
        \emph{standard} remedy is to drop one dummy (reference encoding) or equivalently to
        drop the intercept, restoring full column rank of $\mathbf{X}$. Adding a ridge
        penalty does numerically stabilise $\mathbf{X}^\top \mathbf{X} + \lambda I$, but
        the individual ridge-coefficients of the $K$ dummies are not identified in the
        usual structural sense --- they depend on $\lambda$ and on the (arbitrary) choice
        of intercept column, so this is not the canonical remedy.
  \item \subpoints{1} \textbf{True}. The variance of the joint contribution
        $\mathrm{Var}(\hat\beta_1 X_1 + \hat\beta_2 X_2)$ includes the (large, \emph{negative})
        covariance term that cancels much of the per-coefficient variance inflation; this is why
        prediction can remain stable under near-collinearity even when individual coefficients are
        not, and why the joint test on $\{X_1, X_2\}$ can be highly significant while the marginal
        $t$-tests are not.
\end{enumerate}
\emph{Grading: 1\,P per statement. (iii) is the subtle one and is the deep reason ``$p$-values
of individually collinear predictors are misleading'' --- credit for any answer that names the
covariance cancellation.}

\subsection*{f) Boosting flavors: AdaBoost, gradient boosting, XGBoost \subpoints{3}}
\sol{}
\begin{enumerate}[(i)]
  \item \subpoints{1} \textbf{True}. Each tree contributes $\nu\cdot\hat T_m(x)$, so halving $\nu$
        approximately doubles the $M$ needed to reach the same cumulative ensemble; $(M,\nu)$
        are jointly tuned because they trade off against each other.
  \item \subpoints{1} \textbf{False}. The prof's canonical range is $d \in \{1,\dots,8\}$,
        with stumps ($d=1$) as a routine default; deep trees ($d \in \{10,\dots,50\}$) defeat
        the ``many weak learners'' principle by making each base learner already low-bias and
        high-variance, which boosting cannot easily unwind. Using stumps is the textbook
        recommendation, not a defeat of the bias--variance logic.
  \item \subpoints{1} \textbf{True}. All three claims are lecture-verbatim XGBoost differences:
        (a) leaf-weight penalty $\gamma|T| + \tfrac{1}{2}\lambda\sum_j w_j^2$ in the split
        objective, (b) Newton-style split scoring using both first and second derivatives of
        the loss, (c) row + column subsampling for variance reduction and speed.
\end{enumerate}

\subsection*{g) Logistic regression with categorical $\times$ continuous interaction \subpoints{3}}
\sol{}
\begin{enumerate}[(i)]
  \item \subpoints{1} The odds factor for a unit increase in $x$ depends on the group:
        \begin{itemize}
          \item Group A (reference): $\exp(0.20) \approx \boxed{1.22}$,
          \item Group B: $\exp(0.20 - 0.05) = \exp(0.15) \approx \boxed{1.16}$,
          \item Group C: $\exp(0.20 - 0.15) = \exp(0.05) \approx \boxed{1.05}$.
        \end{itemize}
        \emph{Grading: 0.33\,P each. Common slip: using $\exp(\hat\beta_x) = e^{0.20}$ for all
        three groups (forgetting the interaction).}

  \item \subpoints{1} For group C, $x = 4$:
        $$\hat\eta = -1.50 + 0.20\cdot 4 + 1.20 + (-0.15)\cdot 4 = -1.50 + 0.80 + 1.20 - 0.60 = -0.10.$$
        $$\hat p = \sigma(-0.10) = \frac{1}{1 + e^{0.10}} \approx \boxed{0.475}.$$
        \emph{Grading: 0.5\,P for $\hat\eta$, 0.5\,P for $\hat p$.}

  \item \subpoints{1} \textbf{False}. $\hat\beta_j$ in logistic regression is the change in
        \emph{log-odds} per unit increase in $x_j$ --- equivalently, $\exp(\hat\beta_j)$ is the
        \emph{odds-ratio}. The change in probability is nonlinear in $x_j$ and depends on the
        current value of $\hat p$ (it is largest near $\hat p = 0.5$ and shrinks toward $0$ at
        either extreme).
\end{enumerate}

\subsection*{h) Neural-network regularization \subpoints{3}}
\sol{}
\begin{enumerate}[(i)]
  \item \subpoints{1} With $C = 5$ and $\varepsilon = 0.10$ the off-target mass per class is
        $\varepsilon/(C-1) = 0.10/4 = 0.025$; the on-target mass is $1 - \varepsilon = 0.90$.
        The original one-hot $(0,0,1,0,0)$ becomes
        $$\boxed{(0.025,\; 0.025,\; 0.90,\; 0.025,\; 0.025)}.$$
        \emph{Grading: 1\,P for the full vector, 0.5\,P if only the on-target entry is right.}

  \item \subpoints{1} (A) \textbf{True}. Dropout masks units randomly at train time and is
        switched off at test time; weights are typically rescaled (test-time scaling) or
        activations are scaled at train time (inverted dropout) so that the expected unit
        input matches between regimes.
        (B) \textbf{False}. Early stopping monitors a held-out \emph{validation} set, not the
        training loss. Stopping at the training-loss minimum essentially never stops (training
        loss usually decreases monotonically with more epochs).
        \emph{Grading: 0.5\,P each.}

  \item \subpoints{1} \textbf{False}. The course-recommended dropout range was $0.2$--$0.5$;
        $0.5$ on hidden layers (and even higher on the input layer) is a standard, well-tested
        setting and the prof did not flag it as too aggressive.
        \emph{Grading: 1\,P. Accept ``False'' with any reasonable counter-claim citing the course range.}
\end{enumerate}

\subsection*{i) PCA: variance, scree, score \subpoints{3}}
\sol{}
\begin{enumerate}[(i)]
  \item \subpoints{1} For \emph{standardized} variables each has unit variance, so the total
        variance equals $p = 5$ (and indeed $\sum_j \lambda_j = 2.10 + 1.30 + 0.80 + 0.50 + 0.30 = 5.00$,
        as expected).
        $$\text{PVE}_{1{:}3} = \frac{2.10 + 1.30 + 0.80}{5.00} = \frac{4.20}{5.00} = \boxed{0.84}.$$
        \emph{Grading: 0.5\,P for the total, 0.5\,P for the proportion.}

  \item \subpoints{1} The PC1 score is the inner product $\phi_1^\top x^*$:
        $$z_1^* = 0.55(1.0) + (-0.40)(2.0) + 0.45(-1.0) + 0.50(0.5) + (-0.30)(-0.5)$$
        $$= 0.55 - 0.80 - 0.45 + 0.25 + 0.15 = \boxed{-0.30}.$$
        \emph{Grading: 1\,P. Sign-flip error (giving $+0.30$) gets 0.5\,P if the arithmetic is otherwise clean.}

  \item \subpoints{1} (A) \textbf{True}. Check: $0.55^2 + 0.40^2 + 0.45^2 + 0.50^2 + 0.30^2 =
        0.3025 + 0.16 + 0.2025 + 0.25 + 0.09 = 1.005 \approx 1$ (round-off).
        (B) \textbf{False}. PCA is \emph{not} invariant under rescaling; without standardization the
        variables with the largest raw variance dominate the first few PCs. This is exactly why we
        standardize before PCA.
        (C) \textbf{True}. The $j$th eigenvalue of the sample covariance / correlation matrix is the
        empirical variance of the $j$th score $z_j = \phi_j^\top x$.
        \emph{Grading: 0.33\,P each.}
\end{enumerate}

\vspace{0.6em}
\hrule
\vspace{1em}

%=========================================================
\section*{Problem 3 \points{16} --- Theory, math, pseudocode}
%=========================================================

\subsection*{a) LDA decision boundary and the LDA $\to$ QDA bridge \subpoints{8}}
\sol{}

\begin{enumerate}[(i)]
  \item \subpoints{3} The denominator of Bayes' rule does not depend on $k$, so $\arg\max_k
        \Pr(Y = k \mid X = x) = \arg\max_k \pi_k f_k(x) = \arg\max_k \log[\pi_k f_k(x)]$.
        Take the logarithm:
        $$\log[\pi_k f_k(x)] \;=\; \log \pi_k \;-\; \tfrac{p}{2}\log(2\pi) \;-\; \tfrac{1}{2}\log|\boldsymbol\Sigma|
                                \;-\; \tfrac{1}{2}(x - \boldsymbol\mu_k)^\top \boldsymbol\Sigma^{-1}(x - \boldsymbol\mu_k).$$
        Expand the quadratic:
        $$-\tfrac{1}{2}(x - \boldsymbol\mu_k)^\top \boldsymbol\Sigma^{-1}(x - \boldsymbol\mu_k)
          = -\tfrac{1}{2} x^\top \boldsymbol\Sigma^{-1} x
            + x^\top \boldsymbol\Sigma^{-1} \boldsymbol\mu_k
            - \tfrac{1}{2}\boldsymbol\mu_k^\top \boldsymbol\Sigma^{-1} \boldsymbol\mu_k.$$
        The terms $-\tfrac{p}{2}\log(2\pi)$, $-\tfrac{1}{2}\log|\boldsymbol\Sigma|$, and
        $-\tfrac{1}{2}x^\top \boldsymbol\Sigma^{-1} x$ do \emph{not} depend on the class $k$
        (the first two are constants; the third depends on $x$ alone, which is fixed when we
        do the $\arg\max_k$). Dropping them preserves $\arg\max_k$, leaving exactly
        $$\delta_k(x) \;=\; x^\top \boldsymbol\Sigma^{-1} \boldsymbol\mu_k
                          \;-\; \tfrac{1}{2}\boldsymbol\mu_k^\top \boldsymbol\Sigma^{-1} \boldsymbol\mu_k
                          \;+\; \log \pi_k. \quad\square$$
        \emph{Grading: 3\,P total. 1\,P for ``log of $\pi_k f_k(x)$, denominator dropped,'' 1\,P for
        the quadratic expansion, 1\,P for identifying and dropping the $k$-independent terms.}

  \item \subpoints{2} Setting $\delta_A(x) = \delta_B(x)$:
        $$x^\top \boldsymbol\Sigma^{-1} \boldsymbol\mu_A - \tfrac{1}{2}\boldsymbol\mu_A^\top \boldsymbol\Sigma^{-1} \boldsymbol\mu_A + \log \pi_A
          = x^\top \boldsymbol\Sigma^{-1} \boldsymbol\mu_B - \tfrac{1}{2}\boldsymbol\mu_B^\top \boldsymbol\Sigma^{-1} \boldsymbol\mu_B + \log \pi_B.$$
        Rearrange:
        $$x^\top \boldsymbol\Sigma^{-1} (\boldsymbol\mu_A - \boldsymbol\mu_B)
          \;=\; \tfrac{1}{2}\bigl(\boldsymbol\mu_A^\top \boldsymbol\Sigma^{-1} \boldsymbol\mu_A
                                    - \boldsymbol\mu_B^\top \boldsymbol\Sigma^{-1} \boldsymbol\mu_B\bigr)
                \;-\; \log\frac{\pi_A}{\pi_B}.$$
        This is exactly the claimed equation (the left side is linear in $x$, the right side is
        a scalar constant; hence the boundary is an affine hyperplane). \emph{The boundary is
        linear in $x$ because the quadratic $x^\top \boldsymbol\Sigma^{-1} x$ cancelled when we
        subtracted $\delta_A - \delta_B$ --- this cancellation is possible precisely because both
        classes share the same $\boldsymbol\Sigma$.}
        \emph{Grading: 1\,P for the algebra, 1\,P for the one-sentence reason (shared $\boldsymbol\Sigma$
        $\Rightarrow$ quadratic term cancels).}

  \item \subpoints{2} With $\boldsymbol\mu_A = (1,0)^\top$, $\boldsymbol\mu_B = (0,3)^\top$,
        $\boldsymbol\Sigma = \mathbf{I}_2$, and $\pi_A = \pi_B$ (so $\log(\pi_A/\pi_B) = 0$):
        \begin{align*}
        \boldsymbol\mu_A - \boldsymbol\mu_B &= (1, -3)^\top, \\
        \boldsymbol\mu_A^\top \boldsymbol\mu_A &= 1^2 + 0^2 = 1, \\
        \boldsymbol\mu_B^\top \boldsymbol\mu_B &= 0^2 + 3^2 = 9.
        \end{align*}
        Plug in:
        $$(1, -3) \, x = \tfrac{1}{2}(1 - 9) - 0 = -4
          \quad\Longleftrightarrow\quad \boxed{x_1 - 3 x_2 = -4}.$$
        \emph{Grading: 1\,P for the substitutions, 1\,P for the final equation. Sign-flip in
        $\boldsymbol\mu_A - \boldsymbol\mu_B$ is a common slip; deduct 0.5\,P.}

  \item \subpoints{1} For QDA, $\boldsymbol\Sigma_k$ depends on $k$, so the quadratic term
        $-\tfrac{1}{2} x^\top \boldsymbol\Sigma_k^{-1} x$ in $\delta_k(x)$ no longer cancels
        between classes when one subtracts $\delta_A - \delta_B$. The resulting boundary
        $\delta_A(x) = \delta_B(x)$ therefore contains a non-zero $x^\top(\boldsymbol\Sigma_A^{-1}
        - \boldsymbol\Sigma_B^{-1})x$ term and is quadratic in $x$.
        \emph{Grading: 1\,P. Accept any answer that names the $x^\top \boldsymbol\Sigma^{-1} x$
        term as the surviving culprit.}
\end{enumerate}

\subsection*{b) Pseudocode: $k$-fold CV with the one-SE rule \subpoints{4}}
\sol{}

\begin{enumerate}[(i)]
  \item \subpoints{3}
\begin{verbatim}
Input: dataset (X, y) of size n; grid Lambda = {lambda_1, ..., lambda_T}
       ordered so larger lambda = simpler model; number of folds K = 10.

1. Randomly permute the indices {1, ..., n}; partition into K disjoint
   folds F_1, ..., F_K of (nearly) equal size.

2. For each lambda in Lambda:
     For each fold k = 1, ..., K:
        train_k = (X, y) with rows in F_k removed
        Fit model M_{lambda, k} on train_k
        MSE_{lambda, k} = mean squared error of M_{lambda, k} on F_k
     CV(lambda)  = (1/K) * sum_k MSE_{lambda, k}
     SE(lambda)  = sd(MSE_{lambda, 1}, ..., MSE_{lambda, K}) / sqrt(K)

3. lambda_min  = argmin_lambda CV(lambda)
   bound       = CV(lambda_min) + SE(lambda_min)
   lambda_1SE  = max { lambda in Lambda : CV(lambda) <= bound }
                  (largest = simplest model whose CV is within the SE band)

4. Refit model on the FULL dataset (X, y) at lambda = lambda_1SE.

Return: the refitted model and lambda_1SE.
\end{verbatim}
        \emph{Grading: 3\,P. 0.5\,P for the random partition, 1\,P for the nested
        ``for lambda / for fold'' structure, 0.5\,P for the SE formula, 0.5\,P for the
        one-SE choice as $\max\{\lambda : \mathrm{CV}(\lambda) \le \text{bound}\}$ (the
        \emph{largest} qualifier is essential), 0.5\,P for the final refit on the full data.}

  \item \subpoints{1} The CV curve is itself estimated with noise, so within an SE band of the
        minimum the differences are statistically indistinguishable; the one-SE rule breaks
        the tie in favour of the \emph{simpler} model, which (i) generalizes more reliably out
        of sample and (ii) is more interpretable, at essentially no test-error cost.
        \emph{Grading: 1\,P. Accept ``protect against noise in the CV estimate by picking the
        simplest model in the noise band.''}
\end{enumerate}

\subsection*{c) Backpropagation: sigmoid + binary cross-entropy \subpoints{4}}
\sol{}

\begin{enumerate}[(i)]
  \item \subpoints{2} Step 1: $\partial L / \partial \hat p$. Differentiate
        $L = -[y\log\hat p + (1-y)\log(1-\hat p)]$:
        $$\frac{\partial L}{\partial \hat p}
          \;=\; -\frac{y}{\hat p} + \frac{1-y}{1-\hat p}
          \;=\; \frac{-y(1-\hat p) + (1-y)\hat p}{\hat p(1-\hat p)}
          \;=\; \frac{\hat p - y}{\hat p(1 - \hat p)}.$$
        Step 2: $\partial \hat p / \partial z_2$. Since $\hat p = \sigma(z_2)$,
        $$\frac{\partial \hat p}{\partial z_2} = \sigma(z_2)(1 - \sigma(z_2)) = \hat p(1 - \hat p).$$
        Chain together:
        $$\frac{\partial L}{\partial z_2}
          = \frac{\hat p - y}{\hat p(1-\hat p)} \cdot \hat p(1-\hat p)
          = \boxed{\hat p - y}. \quad\square$$
        The denominator from the cross-entropy derivative cancels exactly with the sigmoid
        derivative --- this is why sigmoid + cross-entropy is the standard pairing (no
        vanishing-gradient problem at saturated outputs, unlike sigmoid + squared error).
        \emph{Grading: 2\,P. 0.5\,P for $\partial L/\partial \hat p$, 0.5\,P for
        $\partial \hat p / \partial z_2$, 1\,P for showing the cancellation.}

  \item \subpoints{1} $z_2 = w_2 h + b_2$, so $\partial z_2 / \partial w_2 = h$. Chain:
        $$\frac{\partial L}{\partial w_2} \;=\; \frac{\partial L}{\partial z_2}\cdot \frac{\partial z_2}{\partial w_2}
          \;=\; \boxed{(\hat p - y)\cdot h}.$$
        \emph{Grading: 1\,P.}

  \item \subpoints{1} Continue the chain through the hidden layer:
        $$\frac{\partial L}{\partial w_1}
          \;=\; \underbrace{(\hat p - y)}_{\partial L/\partial z_2}\cdot
                \underbrace{w_2}_{\partial z_2/\partial h}\cdot
                \underbrace{\mathbb{1}[z_1 > 0]}_{\partial h/\partial z_1}\cdot
                \underbrace{x}_{\partial z_1/\partial w_1}
          \;=\; \boxed{(\hat p - y)\, w_2\, \mathbb{1}[z_1 > 0]\, x}.$$
        \emph{Grading: 1\,P. Common slip: forgetting the ReLU indicator $\mathbb{1}[z_1 > 0]$
        (deduct 0.5\,P).}
\end{enumerate}

\vspace{0.6em}
\hrule
\vspace{1em}

%=========================================================
\section*{Problem 4 \points{20} --- Data analysis: diabetes progression}
%=========================================================

\subsection*{a) OLS with quadratic, interaction, collinear pair \subpoints{7}}
\sol{}
\begin{enumerate}[(i)]
  \item \subpoints{1} The model has $10$ estimated coefficients: intercept, \texttt{age},
        \texttt{sex}, \texttt{bmi}, $I(\texttt{bmi}^2)$, \texttt{bp}, \texttt{s1}, \texttt{s2},
        \texttt{s5}, and \texttt{bmi:sex}. Residual d.f.\ $= n_\text{train} - 10 = 350 - 10 = \boxed{340}$,
        consistent with the printed output.
        \emph{Grading: 1\,P (0.5\,P for the count, 0.5\,P for the d.f.\ check).}

  \item \subpoints{2} \textbf{Near-multicollinearity} between \texttt{s1} (total cholesterol) and
        \texttt{s2} (LDL cholesterol). The two pieces of evidence are: (a) the problem statement
        gives $\mathrm{cor}(\texttt{s1}, \texttt{s2}) \approx 0.97$ on the training set, and (b)
        the SEs of $\hat\beta_{\texttt{s1}}$ and $\hat\beta_{\texttt{s2}}$ are $18.0$ and $17.5$,
        roughly $4$--$6\times$ larger than those of the other (mutually weakly correlated)
        predictors. With $\mathbf{X}^\top \mathbf{X}$ near-singular along the (\texttt{s1, s2})
        direction, $\mathrm{Var}(\hat\beta_j) = \sigma^2 [(\mathbf{X}^\top \mathbf{X})^{-1}]_{jj}$
        explodes for those two coordinates while predictions and the joint contribution remain
        well-identified.
        \emph{Grading: 2\,P. 1\,P for naming collinearity, 1\,P for citing both the correlation and
        the inflated SEs.}

  \item \subpoints{1} Each $p$-value tests $H_0: \beta_j = 0$ \emph{conditional on all other
        predictors, including the other near-collinear partner}, so the test is asking ``can
        \texttt{s1} be set to zero \emph{given} that \texttt{s2} is already in the model?'' (and
        vice versa) --- both answer yes individually, even though dropping \emph{both} typically
        worsens fit substantially. The correct procedure is an $F$-test on the joint
        contribution of $\{\texttt{s1}, \texttt{s2}\}$, or to drop just one (or combine, e.g.\
        via PCA / a sum).
        \emph{Grading: 1\,P. Accept any answer that names the conditional nature of the $t$-tests
        under collinearity.}

  \item \subpoints{2} For each patient, all other (standardized) predictors are at zero, so only
        the intercept, \texttt{sex}, \texttt{bmi}, $I(\texttt{bmi}^2)$, and \texttt{bmi:sex}
        terms contribute.

        Patient A (\texttt{sex}~$= 0$, $\texttt{bmi} = 1$):
        $$\widehat{\texttt{prog}}_A = 152.5 + (-10.2)(0) + 24.0(1) + 3.8(1)^2 + (-9.0)(1)(0)
          = 152.5 + 24.0 + 3.8 = \boxed{180.3}.$$

        Patient B (\texttt{sex}~$= 1$, $\texttt{bmi} = 1$):
        $$\widehat{\texttt{prog}}_B = 152.5 + (-10.2)(1) + 24.0(1) + 3.8(1)^2 + (-9.0)(1)(1)$$
        $$= 152.5 - 10.2 + 24.0 + 3.8 - 9.0 = \boxed{161.1}.$$
        The interaction $\hat\beta_{\texttt{bmi:sex}} = -9.0$ implies the slope of \texttt{prog}
        in \texttt{bmi} is $9.0$ units smaller for males than for females.
        \emph{Grading: 1\,P each. Forgetting the \texttt{bmi:sex} term for patient B is a hard
        miss (deduct 0.5\,P).}

  \item \subpoints{1} \emph{Convention:} standardized $\texttt{bmi}$ has mean $0$ and SD $1$,
        so the values $+1$ and $+2$ correspond to one and two SDs above the training mean of
        raw BMI. For a female (\texttt{sex}~$=0$), going from $\texttt{bmi} = 1$ to $\texttt{bmi} = 2$:
        \begin{align*}
        \Delta_\text{linear}    &= 24.0 \cdot (2 - 1) = 24.0, \\
        \Delta_\text{quadratic} &= 3.8 \cdot (2^2 - 1^2) = 3.8 \cdot 3 = 11.4, \\
        \Delta_\text{interaction} &= -9.0 \cdot 0 \cdot (2 - 1) = 0, \\
        \Delta_\text{total}     &= 24.0 + 11.4 = \boxed{35.4}.
        \end{align*}
        \emph{Grading: 1\,P. The trap is to forget the quadratic contribution and answer
        $24.0$; deduct 0.5\,P.}
\end{enumerate}

\subsection*{b) Lasso with 10-fold CV \subpoints{4}}
\sol{}
\begin{enumerate}[(i)]
  \item \subpoints{1} With $p = 9$ predictors and a standardized design (no intercept column in
        $\mathbf{X}$), the lasso objective is
        $$\min_{\beta_0, \boldsymbol\beta} \;\tfrac{1}{2}\bigl\|\mathbf{y} - \beta_0 \mathbf{1}_n - \mathbf{X}\boldsymbol\beta\bigr\|_2^2
          \;+\; \lambda \sum_{j=1}^{p} \lvert \beta_j \rvert,$$
        where the intercept $\beta_0$ is \emph{not} penalized (and is typically eliminated by
        centring the response) and the columns of $\mathbf{X}$ are standardized to make the
        $L^1$ penalty scale-invariant.
        \emph{Grading: 1\,P. Penalizing the intercept is the canonical wrong answer; deduct fully.}

  \item \subpoints{1} The dummy-coding scheme determines which level becomes the ``zero''
        reference, and the lasso penalty is applied to the contrast coefficients, not to a
        symmetric group penalty (unless one uses group lasso); so different reference levels
        give different effective penalties on the same categorical predictor, and which dummy
        gets shrunk to zero is a function of the coding choice rather than of the data. The
        practical consequence: a dummy shrunk to zero \emph{means the corresponding level
        cannot be distinguished from the reference}, not that the predictor is irrelevant.
        \emph{Grading: 1\,P.}

  \item \subpoints{1} Lasso reduces the \emph{variance} of the OLS predictions at the cost of
        introducing some bias. The near-collinear pair (\texttt{s1, s2}) inflates the variance
        of the OLS predictions through their coefficient variances; lasso resolves the
        ambiguity by shrinking one of them (here: 7 nonzeros out of 9), which substantially
        lowers prediction variance, and that variance reduction outweighs the bias introduced
        by shrinkage.
        \emph{Grading: 1\,P. Accept ``variance reduction at the cost of some bias, exploiting
        the collinear pair where OLS variance was inflated.''}

  \item \subpoints{1} \textbf{One-SE rule:} from the CV curve at $\hat\lambda_{\min}$, compute
        $\widehat{\mathrm{SE}}(\mathrm{CV}_K)$ and pick the \emph{simplest} model (largest
        $\lambda$, fewest nonzero coefficients) whose CV is within one SE of the minimum.
        The practitioner's preference is reasonable because (i) the differences within the SE
        band are statistically indistinguishable, and (ii) a model with $5$ predictors instead
        of $7$ is more stable across resamples and easier to interpret --- both core arguments
        for the rule.
        \emph{Grading: 1\,P (0.5 for the rule, 0.5 for a reasonable justification).}
\end{enumerate}

\subsection*{c) Gradient boosting --- learning-rate / depth tradeoff \subpoints{5}}
\sol{}
\begin{enumerate}[(i)]
  \item \subpoints{2} Squared-error gradient boosting:
\begin{verbatim}
Initialize  f^{(0)}(x) = mean(y)              [the constant minimizing SSE]
For m = 1, ..., M:
    r_i = y_i - f^{(m-1)}(x_i)               [negative gradient = residual]
    Fit a regression tree h_tilde_m to (X, r) using subroutine T(.)
    Update  f^{(m)}(x) = f^{(m-1)}(x) + nu * h_tilde_m(x)
Return f^{(M)}
\end{verbatim}
        \emph{Grading: 2\,P. 0.5\,P initialisation, 0.5\,P fitting to residuals (= negative
        gradient under SSE), 0.5\,P the $\nu$-scaled update, 0.5\,P the final $f^{(M)}$.}

  \item \subpoints{1} Halving $\nu$ approximately \emph{doubles} the required $M^*$ (here:
        $\nu = 0.10 \to M^* = 700$, $\nu = 0.05 \to M^* = 1{,}400$, $\nu = 0.01 \to M^* = 6{,}500$;
        the last factor is $\approx 9$ rather than $5$ because smaller $\nu$ also reduces effective
        bias more gradually). Smaller $\nu$ slows the rate at which the ensemble fits, lowering
        variance per added tree at the cost of needing more trees to reach the same bias.
        \emph{Grading: 1\,P. Accept ``$M^* \cdot \nu \approx$ constant.''}

  \item \subpoints{1} Deep trees are themselves low-bias / high-variance learners; combined
        with \emph{small} $\nu$ the ensemble takes many slow steps each adding lots of variance,
        and combined with \emph{large} $\nu$ each big step over-fits the local residual --- in
        either case boosting can't undo the variance contributed by a non-weak base learner,
        whereas shallow trees keep each step weak so that $\nu$ alone controls the bias--variance
        budget.
        \emph{Grading: 1\,P. Any answer naming the ``deep tree = strong base learner =
        boosting can't compensate'' argument.}

  \item \subpoints{1} \textbf{Unsound} for gradient boosting --- $M$ is a genuine tuning
        parameter and too-large $M$ \emph{does} overfit (test error eventually rises). Random
        forests behave differently: more trees only reduces variance, the ensemble does not
        overfit in $B$, so ``set $B = 500$ to be safe'' is fine for RF but not for boosting.
        \emph{Grading: 1\,P.}
\end{enumerate}

\subsection*{d) GAM and random forest \subpoints{4}}
\sol{}
\begin{enumerate}[(i)]
  \item \subpoints{1} A cubic regression spline with $K = 2$ interior knots has
        $K + d + 1 = 2 + 3 + 1 = 6$ basis functions; removing the global intercept leaves
        $\boxed{5}$ degrees of freedom for the $\mathrm{bs}(\texttt{bp},\dots)$ term.
        \emph{Grading: 1\,P.}

  \item \subpoints{1} Sum up:
        \begin{itemize}
          \item intercept: $1$;
          \item $s(\texttt{age}, \mathrm{df}=4)$: $4$;
          \item $\texttt{sex}$ (single dummy): $1$;
          \item $s(\texttt{bmi}, \mathrm{df}=5)$: $5$;
          \item $\mathrm{bs}(\texttt{bp}, \dots)$ (from (i)): $5$;
          \item $\texttt{s1}, \texttt{s2}, \texttt{s5}$ (linear): $1 + 1 + 1 = 3$.
        \end{itemize}
        Total: $1 + 4 + 1 + 5 + 5 + 3 = \boxed{19}$ degrees of freedom.
        \emph{Grading: 1\,P. Off-by-one slips on the spline conventions are common; accept
        any answer between $18$ and $20$ with correct accounting per term.}

  \item \subpoints{1} For regression random forests, the standard default is $\texttt{mtry} = p/3
        = 9/3 = \boxed{3}$. One picks $\texttt{mtry} < p$ to \emph{decorrelate} the bagged trees:
        bagging only reduces variance when the trees are weakly correlated, and forcing each
        split to consider only a random subset of features prevents the same one or two strong
        predictors (here: \texttt{bmi}, \texttt{s5}, \texttt{bp}) from dominating every tree.
        \emph{Grading: 1\,P.}

  \item \subpoints{1} The random-forest variable-importance plot reports a \emph{magnitude}
        of effect (e.g.\ mean decrease in OOB MSE) but \emph{not a sign / direction}: it tells
        you which predictors matter but not whether higher \texttt{bmi} is associated with
        higher or lower disease progression. An OLS coefficient does both (and supplies an SE
        for inference).
        \emph{Grading: 1\,P. Accept ``no sign'' or ``no functional form'' or ``no confidence
        interval / inference.''}
\end{enumerate}

\vspace{0.6em}
\hrule
\vspace{1em}

%=========================================================
\section*{Problem 5 \points{26} --- Data analysis: South African heart disease}
%=========================================================

\subsection*{a) Logistic regression with a continuous-by-binary interaction \subpoints{8}}
\sol{}
\begin{enumerate}[(i)]
  \item \subpoints{2} \emph{Encoding:} \texttt{famhist}~$=1$ means family history, $0$ means
        no family history (as given). Each additional mmol/L of \texttt{ldl} changes the
        log-odds by $\hat\beta_{\texttt{ldl}} + \hat\beta_{\texttt{ldl:famhist}}\cdot
        \texttt{famhist}$:
        \begin{itemize}
          \item \texttt{famhist}~$=0$: factor $= \exp(0.18) \approx \boxed{1.197}$.
          \item \texttt{famhist}~$=1$: factor $= \exp(0.18 + 0.15) = \exp(0.33) \approx \boxed{1.391}$.
        \end{itemize}
        \emph{Grading: 1\,P each. Forgetting the interaction in the second case (answering
        $1.197$ twice) is a hard miss.}

  \item \subpoints{1} The interaction $\hat\beta_{\texttt{ldl:famhist}} = 0.15$ is non-zero,
        so the slope of the log-odds in \texttt{ldl} \emph{depends on} \texttt{famhist} and the
        ``effect of LDL'' is different in the two groups (specifically: LDL is a stronger
        risk factor among patients with family history). $\hat\beta_{\texttt{ldl}} = 0.18$
        alone reports only the slope in the reference group (\texttt{famhist}~$=0$) and is
        therefore not a useful summary of ``the effect of LDL'' in this fitted model.
        \emph{Grading: 1\,P.}

  \item \subpoints{3} Compute $\hat\eta$ term by term:
        \begin{align*}
        \hat\eta = &\;-5.50  \quad\text{(intercept)} \\
                  &+ 0.045\cdot 55 = +2.475   \quad\text{(\texttt{age})} \\
                  &+ 0.18\cdot 5   = +0.900   \quad\text{(\texttt{ldl})} \\
                  &+ 0.85\cdot 1   = +0.850   \quad\text{(\texttt{famhist})} \\
                  &+ 0.080\cdot 4  = +0.320   \quad\text{(\texttt{tobacco})} \\
                  &+ 0.035\cdot 55 = +1.925   \quad\text{(\texttt{typea})} \\
                  &+ 0.020\cdot 25 = +0.500   \quad\text{(\texttt{adiposity})} \\
                  &-0.050\cdot 27 = -1.350   \quad\text{(\texttt{obesity})} \\
                  &+ 0.15\cdot 5\cdot 1 = +0.750   \quad\text{(\texttt{ldl:famhist})}
        \end{align*}
        Cumulative sum: $-5.50 + 2.475 - 3.025;\ +0.900 \to -2.125;\ +0.850 \to -1.275;\ +0.320 \to -0.955;\
        +1.925 \to +0.970;\ +0.500 \to +1.470;\ -1.350 \to +0.120;\ +0.750 \to$
        $$\boxed{\hat\eta = 0.870}.$$
        $$\hat p \;=\; \sigma(0.870) \;=\; \frac{1}{1 + e^{-0.870}} \;\approx\; \boxed{0.705}.$$
        \emph{Grading: 3\,P. 1.5\,P for $\hat\eta$ (deduct 0.5\,P per arithmetic slip up to
        a floor of $1$\,P if the structure is right), 1\,P for the sigmoid evaluation,
        0.5\,P for correctly including the interaction term.}

  \item \subpoints{1} $\hat p = 0.705 > 0.5$, so the patient \emph{would} be flagged as CHD at the
        default threshold. However $0.5$ is rarely the right operating point for medical
        screening: the base rate is $\approx 35\%$, false negatives are clinically costly, and the
        threshold should be lowered (perhaps to $0.3$) to raise sensitivity at the cost of some
        false positives.
        \emph{Grading: 1\,P. Accept ``yes, flagged, but $0.5$ is too high for screening with
        asymmetric costs.''}

  \item \subpoints{1} The model is an observational regression and reports \emph{conditional
        associations, not causal effects}; the negative sign on \texttt{obesity} could be a
        product of confounding (e.g.\ obesity is correlated with \texttt{adiposity}, \texttt{bmi},
        and \texttt{bp}, and ``controlling for'' those other predictors makes the residual
        contribution of obesity to CHD risk hard to interpret). The prof's verbatim line:
        ``these are fancy correlations, not causal claims.''
        \emph{Grading: 1\,P. Accept any answer naming confounding / observational design /
        ``correlation is not causation.''}
\end{enumerate}

\subsection*{b) AdaBoost: deriving the classifier weight $\alpha_m$ \subpoints{6}}
\sol{}
\begin{enumerate}[(i)]
  \item \subpoints{3} Split the sum by correct vs.\ incorrect prediction of $G$:
        \begin{align*}
        Q(\alpha, G) &= \sum_i w_i^{(m)} \exp(-\alpha\, y_i G(x_i)) \\
                     &= \underbrace{\sum_{i:\, y_i = G(x_i)} w_i^{(m)} e^{-\alpha}}_{W_C \, e^{-\alpha}}
                      \;+\; \underbrace{\sum_{i:\, y_i \neq G(x_i)} w_i^{(m)} e^{+\alpha}}_{W_W \, e^{+\alpha}},
        \end{align*}
        where $W_C = \sum_{i:\,\text{correct}} w_i^{(m)}$ and $W_W = \sum_{i:\,\text{wrong}} w_i^{(m)}$.
        Differentiate with respect to $\alpha$ and set to $0$:
        $$\frac{\partial Q}{\partial \alpha} = -W_C e^{-\alpha} + W_W e^{+\alpha} = 0
          \quad\Longrightarrow\quad e^{2\alpha} = \frac{W_C}{W_W}
          \quad\Longrightarrow\quad \alpha^* = \tfrac{1}{2}\log \frac{W_C}{W_W}.$$
        Define the weighted error $\mathrm{err}_m = W_W / (W_C + W_W)$, so
        $W_C / W_W = (1 - \mathrm{err}_m)/\mathrm{err}_m$, giving
        $$\boxed{\alpha^* = \tfrac{1}{2}\log \frac{1 - \mathrm{err}_m}{\mathrm{err}_m}}. \quad\square$$
        AdaBoost.M1 absorbs the $\tfrac{1}{2}$ into the weight-update rule, so the algorithm uses
        $\alpha_m = \log[(1 - \mathrm{err}_m)/\mathrm{err}_m]$ rather than half of it.
        \emph{Grading: 3\,P. 1\,P for the correct/wrong split, 1\,P for the derivative and
        $e^{2\alpha} = W_C/W_W$, 1\,P for substituting $\mathrm{err}_m$.}

  \item \subpoints{1} For $\mathrm{err}_m = 0.30$:
        $$\alpha_m = \log\frac{1 - 0.30}{0.30} = \log\frac{0.70}{0.30} = \log(2.333\ldots) \approx \boxed{0.85}.$$
        $$\exp(\alpha_m) = \frac{0.70}{0.30} \approx \boxed{2.33}.$$
        \emph{Grading: 0.5\,P each.}

  \item \subpoints{1} If $\mathrm{err}_m > 0.5$, then $\alpha_m = \log[(1-\mathrm{err}_m)/\mathrm{err}_m] < 0$;
        the base learner is worse than random and its sign is flipped in the ensemble vote
        (equivalently, the algorithm trusts $-G_m$). In practice AdaBoost is usually stopped
        or the learner re-selected when this happens.
        \emph{Grading: 1\,P. Accept ``$\alpha_m$ becomes negative; $G_m$ contributes with
        inverted sign.''}

  \item \subpoints{1} A deep tree is already a low-bias / high-variance learner on its own.
        Boosting's strategy is to chain many \emph{weak} learners so that each iteration reduces
        bias by a small, controlled amount; if each base learner is itself high-variance,
        the cumulative variance of the ensemble explodes and the ``many weak learners'' logic
        breaks --- which is why AdaBoost (and most boosting) defaults to stumps or very shallow trees.
        \emph{Grading: 1\,P.}
\end{enumerate}

\subsection*{c) Neural-network classifier with regularization \subpoints{6}}
\sol{}
\begin{enumerate}[(i)]
  \item \subpoints{1} Parameter count layer by layer:
        \begin{itemize}
          \item Hidden layer 1 ($7 \to 20$): $7\cdot 20 = 140$ weights $+ 20$ biases $= 160$.
          \item Hidden layer 2 ($20 \to 10$): $20\cdot 10 = 200$ weights $+ 10$ biases $= 210$.
          \item Output layer ($10 \to 1$): $10\cdot 1 = 10$ weights $+ 1$ bias $= 11$.
        \end{itemize}
        Total: $160 + 210 + 11 = \boxed{381}$ parameters.
        \emph{Grading: 1\,P. Forgetting biases (giving $350$) is a common miss; deduct 0.5\,P.}

  \item \subpoints{1} Pre-activation:
        \begin{align*}
        z = w^\top x + b
          &= 0.10(0.5) + 0.30(1.0) + (-0.20)(-1.0) + 0.50(2.0) \\
          &\quad + 0.05(0.5) + (-0.40)(0.5) + 0.20(-0.5) + (-0.10) \\
          &= 0.05 + 0.30 + 0.20 + 1.00 + 0.025 - 0.20 - 0.10 - 0.10 \\
          &= 1.175.
        \end{align*}
        Since $z > 0$, $h = \mathrm{ReLU}(z) = z \approx \boxed{1.18}$ (and $z \approx \boxed{1.18}$).
        \emph{Grading: 1\,P (0.5\,P for $z$, 0.5\,P for $h$).}

  \item \subpoints{1} Train-test gap of $98\% \to 70\%$ is classic overfitting: with $381$
        parameters and only $\sim 350$ training points, the network has enough capacity to
        memorize the training set, driving training error to near-zero (low \emph{empirical}
        bias) while estimator variance is large and test error is much worse. Regularization
        forces the optimizer toward solutions with smaller capacity / smoother decision
        surfaces, trading a small increase in training error for a substantial decrease in
        test error.
        \emph{Grading: 1\,P.}

  \item \subpoints{1} (a) \textbf{$L^2$ weight decay} adds $\lambda \|\boldsymbol\theta\|_2^2$ to
        the training loss, which shrinks weights toward zero each step and biases the
        optimizer toward smaller-norm (smoother) solutions.
        (b) \textbf{Dropout} randomly masks a fraction $p$ of hidden units at each training
        forward pass, preventing the network from depending on any single feature
        (``ensemble-of-subnetworks'' interpretation); at test time \emph{all} units are kept
        and the weights (or activations) are rescaled by the keep-probability so that the
        expected unit input matches the training-time distribution.
        (c) \textbf{Early stopping} monitors held-out validation loss and halts training at
        the epoch where validation loss first stops improving, returning those weights ---
        which effectively limits the number of optimizer steps and so the effective capacity.
        \emph{Grading: 1\,P total (roughly 0.33\,P each).}

  \item \subpoints{1} \emph{Treating the binary output as the two-class one-hot $(0,1)$ for
        $y = 1$ and applying $\varepsilon/(C-1)$ with $C = 2$:} off-target mass $0.10/1 = 0.10$,
        on-target mass $1 - 0.10 = 0.90$, so the modified target is
        $$\boxed{(0.10,\; 0.90)}.$$
        Label smoothing prevents the network from being forced to drive its softmax / sigmoid
        outputs to exactly $0$ or $1$ (which would require unbounded logits) and acts as a
        regularizer that is also robust to mislabelled training examples --- the gradient
        contribution from a wrong label is bounded rather than infinite.
        \emph{Grading: 1\,P. Accept either symmetric convention $(0.10, 0.90)$ or $(0.05, 0.95)$
        with an explicit statement of which convention is used; the exam's P2(h)(i) used the
        symmetric $\varepsilon/(C-1)$ form which here yields $(0.10, 0.90)$.}

  \item \subpoints{1} The classical bias--variance bound that ``$p \gg n$ must overfit'' is
        false in the over-parameterized regime: SGD's implicit bias toward minimum-norm
        interpolators, combined with regularization (weight decay / dropout / early stopping),
        can produce models that interpolate the training set and still generalize well ---
        the ``benign overfitting'' / double-descent phenomenon. The friend's claim assumes the
        classical regime where it does not apply.
        \emph{Grading: 1\,P.}
\end{enumerate}

\subsection*{d) Comparison of classifiers and class imbalance \subpoints{4}}
\sol{}
\begin{enumerate}[(i)]
  \item \subpoints{2} Logistic regression:
        \begin{itemize}
          \item Sensitivity $= 18 / (18 + 20) = 18/38 \approx \boxed{0.47}$.
          \item Specificity $= 64 / (10 + 64) = 64/74 \approx \boxed{0.86}$.
          \item Error rate $= (20 + 10)/112 = 30/112 \approx \boxed{0.27}$.
        \end{itemize}
        AdaBoost:
        \begin{itemize}
          \item Sensitivity $= 25 / (25 + 13) = 25/38 \approx \boxed{0.66}$.
          \item Specificity $= 60 / (14 + 60) = 60/74 \approx \boxed{0.81}$.
          \item Error rate $= (13 + 14)/112 = 27/112 \approx \boxed{0.24}$.
        \end{itemize}
        \emph{Grading: 0.33\,P per number.}

  \item \subpoints{1} The naive ``no CHD'' classifier misclassifies the $38$ true positives,
        giving error $= 38/112 \approx \boxed{0.34}$. Accuracy alone is a poor discriminator
        here: the naive classifier already has $66\%$ accuracy, AdaBoost has $76\%$ and
        logistic regression $73\%$, all in the same range, and accuracy hides the fact that
        the naive classifier has $0\%$ sensitivity --- it never flags a true CHD patient.
        \emph{Grading: 1\,P.}

  \item \subpoints{1} For a screening task with costly false negatives, the relevant metric
        is \textbf{sensitivity} (recall on the positive class). AdaBoost has sensitivity
        $0.66$ vs.\ logistic regression's $0.47$, so \textbf{deploy AdaBoost}; the cost is a
        modest reduction in specificity ($0.81$ vs.\ $0.86$), which under asymmetric costs is
        the right trade.
        \emph{Grading: 1\,P (accept any answer that names the right classifier and the right
        metric; deduct 0.5\,P for a correct metric with the wrong classifier).}
\end{enumerate}

\subsection*{e) KNN with mixed-unit predictors \subpoints{2}}
\sol{}
\begin{enumerate}[(i)]
  \item \subpoints{1} For a test point $x_0$, compute distances to all training points, pick
        the $K$ closest, and classify $x_0$ by majority vote among those $K$ neighbours'
        labels (with ties broken at random or by a tiebreak rule).
        \emph{Grading: 1\,P.}

  \item \subpoints{1} Euclidean distance on raw predictors lets the variable with the largest
        numerical range (here \texttt{age} in years, with values up to $\sim 60$) dominate
        the distance, swamping the contribution of all other features --- so KNN effectively
        becomes a $1$-D classifier on \texttt{age}. The standard remedy is to
        \emph{standardize} (center and scale to unit SD) each continuous predictor before
        fitting KNN.
        \emph{Grading: 1\,P.}
\end{enumerate}

\vspace{1em}
\hrule
\vspace{0.5em}

\noindent\textbf{End of solution.}\quad Total: $10 + 28 + 16 + 20 + 26 = 100$ points.

\end{document}
